\section{Related Work}
\label{sec:related}

\para{Honesty versus accuracy.}
\mask elicits a model's belief under a neutral prompt and then applies pressure; a statement that contradicts the belief is a lie, and a belief that contradicts the ground truth is an inaccuracy \citep{ren2025mask}.
\citet{pacchiardi2023catch} similarly define a lie as a false output from a model that demonstrably knows the truth.
These works give the behavioral partition we need, but \mask evaluates no detector.
We adopt its belief-elicitation logic and add detectors on top.

\para{Deception probes.}
\citet{goldowskydill2025detecting} train logistic-regression probes on instructed pairs of honest and dishonest personas and report AUROC of 0.96 to 0.999 on role-play, insider-trading, and sandbagging data.
They also note that the probe fires on honest text in deception-related contexts.
\liarsbench evaluates this style of probe on seven lie datasets and four models and finds an average AUROC of 0.60, with scores at or below chance on harm-pressure sets, while an in-distribution probe reaches 0.91 \citep{kretschmar2025liars}.
\citet{natarajan2026one} show that the choice of contrastive prompt pair explains most of the variance in probe performance.
\citet{moustafa2026beyond} find below-chance transfer to harm-pressure sets and no gain from more elaborate probes.
\citet{luikham2026asymmetries} compares instructed and incentive-only lies in a 70B model and finds that fabrication under pressure alone is rare.
\citet{cooney2026did} study a self-report probe on \triviaqa after filtering to questions the model answers correctly, and stress that belief verification is needed to tell a lie from an error.
\citet{hopkins2026finetuned} report that fine-tuned detectors generalize poorly and that many contradiction-based lie labels are overturned on review.
\citet{vonklinski2026stress} argue that probes track confounds such as instruction compliance and response likelihood.
Unlike these works, we score the probe on false answers the model gives without knowing the truth, under the same prompt as the lies.

\para{Truth probes.}
Linear probes can read the truth value of statements from hidden states \citep{burns2022discovering,azaria2023internal,marks2023geometry,burger2024truth}, although they generalize unevenly \citep{levinstein2023still}.
\citet{thormann2026probing} shows that a truth probe flags lies but not true-but-misleading statements, so it tracks falsity and not deceptive intent.
Our question is the mirror image: falsity without deceptive intent.

\para{Hallucination detectors.}
Correctness probes, answer probability, and sample consistency predict whether an answer is wrong \citep{orgad2024llms,kossen2024semantic}.
\citet{cheang2025llms} find that these detectors reach AUROC 0.86 to 0.93 on hallucinations where the model recalls nothing about the subject and only 0.48 to 0.69 where it recalls a wrong association, which suggests that they read knowledge recall.
Their study has no lie condition.
We test the same detector family on knowing lies.

\para{Lies versus hallucinations.}
\citet{huan2025can} contrast lying with hallucination mechanistically.
Two recent preprints approach our question with other tools: one proposes a residual-rank signature that separates knowing lies from fine-tuned ``naive liar'' controls \citep{nyoma2026rift}, and one uses recognition probing to separate ``will not answer'' from ``cannot answer'' \citep{dingeto2026lie}.
We read these two only as abstracts, so we do not claim that our design is unprecedented.
Building on the works above, we cross knowledge with incentive on matched questions and apply both detector families to every cell.
